\chapter{Algoritmos em Strings}

%==============================================================================
\section{Hash de strings}
%==============================================================================

\subsection{Padrão do Problema}
Comparar substrings rapidamente, detectar repetições ou usar strings como estados sem comparar caractere por caractere.

\subsection{Quando usar}
\begin{itemize}[nosep]
  \item Há muitas consultas perguntando se duas substrings são iguais.
  \item É preciso testar padrões, prefixos, sufixos ou repetições em uma string grande.
  \item Uma comparação direta de substrings ficaria cara demais.
  \item Probabilidade pequena de colisão é aceitável, especialmente com hash duplo.
\end{itemize}

\subsection{Explicação do Algoritmo}
Pré-calcule hashes de prefixos e potências da base. O hash de qualquer substring é obtido em $O(1)$ removendo a contribuição do prefixo anterior e ajustando pelo tamanho.

\subsection{Complexidade Temporal}
Pré-processamento $O(n)$ e consulta $O(1)$.

\subsection{Exemplo de Enunciado}
\begin{exemplodiabox}
\textbf{Enunciado:} Dada uma string $s$ e várias consultas $(l_1,r_1,l_2,r_2)$, diga se as duas substrings indicadas são iguais.
\end{exemplodiabox}

\subsection{Código}
\incluircppcaption{algoritmos_avulsos/strings/hash.cpp}{Hash duplo de strings}

%==============================================================================
\section{Função de prefixo (LPS)}
%==============================================================================

\subsection{Padrão do Problema}
Calcular, para cada posição, o maior prefixo próprio da string que também é sufixo do prefixo atual.

\subsection{Quando usar}
\begin{itemize}[nosep]
  \item O problema envolve bordas de uma string.
  \item É preciso preparar a string para KMP.
  \item O enunciado pede períodos, prefixos que também são sufixos ou repetições.
\end{itemize}

\subsection{Explicação do Algoritmo}
A função de prefixo mantém o tamanho da melhor borda até a posição anterior. Quando o próximo caractere não encaixa, volta usando valores já calculados até encontrar uma borda menor válida.

\subsection{Complexidade Temporal}
$O(n)$.

\subsection{Exemplo de Enunciado}
\begin{exemplodiabox}
\textbf{Enunciado:} Para cada prefixo de uma string, calcule o tamanho do maior prefixo próprio que também é sufixo.

\textbf{Exemplo:} Em \texttt{ababaca}, a informação de bordas permite continuar buscas sem voltar no texto.
\end{exemplodiabox}

\subsection{Código}
\incluircppcaption{algoritmos_avulsos/strings/kmp_lps.cpp}{Função de prefixo}

%==============================================================================
\section{KMP}
%==============================================================================

\subsection{Padrão do Problema}
Encontrar todas as ocorrências de um padrão dentro de um texto.

\subsection{Quando usar}
\begin{itemize}[nosep]
  \item É necessário buscar um padrão exato em um texto.
  \item O texto e o padrão podem ser grandes.
  \item Uma busca ingênua $O(nm)$ pode estourar o tempo.
  \item O problema pede todas as posições de ocorrência.
\end{itemize}

\subsection{Explicação do Algoritmo}
KMP usa a função de prefixo do padrão para saber quanto do casamento pode ser reaproveitado após uma falha. Assim, o ponteiro do texto nunca volta.

\subsection{Complexidade Temporal}
$O(n + m)$, com $n$ tamanho do texto e $m$ tamanho do padrão.

\subsection{Exemplo de Enunciado}
\begin{exemplodiabox}
\textbf{Enunciado:} Dado um texto e uma palavra, imprima todos os índices em que a palavra aparece no texto.

\textbf{Exemplo:} Texto \texttt{ababa}, padrão \texttt{aba} $\Rightarrow$ ocorrências em $0$ e $2$.
\end{exemplodiabox}

\subsection{Código}
\incluircppcaption{algoritmos_avulsos/strings/kmp.cpp}{KMP}

%==============================================================================
\section{Rabin-Karp}
%==============================================================================

\subsection{Padrão do Problema}
Buscar padrões por hash, especialmente quando há muitas comparações de substrings de mesmo tamanho.

\subsection{Quando usar}
\begin{itemize}[nosep]
  \item É preciso procurar um padrão em um texto usando rolling hash.
  \item Há muitas substrings candidatas do mesmo tamanho.
  \item O risco de colisão é aceitável ou pode ser reduzido com hash duplo.
  \item O problema combina busca de padrão com outras comparações por hash.
\end{itemize}

\subsection{Explicação do Algoritmo}
Calcule o hash do padrão e hashes de prefixos do texto. Para cada janela do tamanho do padrão, compare o hash da janela com o hash do padrão ajustado pela potência correspondente.

\subsection{Complexidade Temporal}
$O(n + m)$.

\subsection{Exemplo de Enunciado}
\begin{exemplodiabox}
\textbf{Enunciado:} Dado um padrão $p$ e um texto $t$, encontre todas as posições em que $p$ aparece em $t$ usando hashing.
\end{exemplodiabox}

\subsection{Código}
\incluircppcaption{algoritmos_avulsos/strings/rabin_karp.cpp}{Rabin-Karp}

%==============================================================================
\section{Manacher}
%==============================================================================

\subsection{Padrão do Problema}
Encontrar palíndromos em uma string, calculando o maior raio palindrômico em cada centro.

\subsection{Quando usar}
\begin{itemize}[nosep]
  \item O problema pede maior substring palindrômica.
  \item É necessário contar ou consultar palíndromos por centro.
  \item Uma expansão de centro $O(n^2)$ é lenta demais.
  \item A string tem tamanho grande.
\end{itemize}

\subsection{Explicação do Algoritmo}
Manacher reaproveita informações do maior palíndromo conhecido à direita. Para cada centro, inicia com um raio mínimo espelhado e expande apenas quando necessário.

\subsection{Complexidade Temporal}
$O(n)$.

\subsection{Exemplo de Enunciado}
\begin{exemplodiabox}
\textbf{Enunciado:} Dada uma string, encontre o tamanho da maior substring palindrômica.

\textbf{Exemplo:} Em \texttt{babad}, uma resposta possível é \texttt{bab}.
\end{exemplodiabox}

\subsection{Código}
\incluircppcaption{algoritmos_avulsos/strings/manacher.cpp}{Manacher}

%==============================================================================
\section{Trie}
%==============================================================================

\subsection{Padrão do Problema}
Guardar várias palavras e responder consultas por prefixo.

\subsection{Quando usar}
\begin{itemize}[nosep]
  \item O problema envolve dicionário de palavras.
  \item Há consultas de prefixo, inserção ou remoção.
  \item Comparar todas as palavras a cada consulta seria lento.
  \item O alfabeto é pequeno ou controlado, como letras minúsculas.
\end{itemize}

\subsection{Explicação do Algoritmo}
Trie é uma árvore em que cada aresta representa um caractere. Palavras com prefixo comum compartilham o mesmo caminho inicial, permitindo buscar ou inserir em tempo proporcional ao tamanho da palavra.

\subsection{Complexidade Temporal}
$O(L)$ por operação, com $L$ tamanho da palavra consultada.

\subsection{Exemplo de Enunciado}
\begin{exemplodiabox}
\textbf{Enunciado:} Mantenha um conjunto de palavras e responda quantas palavras começam com determinado prefixo.
\end{exemplodiabox}

\subsection{Código}
\incluircppcaption{algoritmos_avulsos/strings/trie.cpp}{Trie}

%==============================================================================
\section{Autômato de sufixos}
%==============================================================================

\subsection{Padrão do Problema}
Representar todas as substrings de uma string em uma estrutura compacta.

\subsection{Quando usar}
\begin{itemize}[nosep]
  \item É preciso contar substrings distintas.
  \item O problema pede maior substring comum entre duas strings.
  \item Há várias consultas sobre substrings de uma string fixa.
  \item Suffix array ou hashing ficariam menos diretos para a operação pedida.
\end{itemize}

\subsection{Explicação do Algoritmo}
O autômato de sufixos mantém estados que representam conjuntos de substrings com mesmos finais possíveis. Ao inserir cada caractere, cria ou clona estados para preservar as transições e links de sufixo.

\subsection{Complexidade Temporal}
Construção $O(n \log \sigma)$ usando \texttt{map}; com array para alfabeto fixo, $O(n)$.

\subsection{Exemplo de Enunciado}
\begin{exemplodiabox}
\textbf{Enunciado:} Dada uma string $s$, conte quantas substrings distintas ela possui.

\textbf{Exemplo:} Para \texttt{ababa}, o autômato permite contar todas as substrings diferentes em tempo linear após a construção.
\end{exemplodiabox}

\subsection{Código}
\incluircppcaption{algoritmos_avulsos/strings/suffix_automaton.cpp}{Autômato de sufixos}
